\subsection{The mixed determinant inequality for the actual forms}
\label{actual-mixed-determinant-cauchy}

The next Fourier input compares the plus and minus residue words.
An exact inverse identity supplies it on the original coefficient
classes, even when the underlying residue form has a kernel.

Let $w$ be an original signed word and let $y_1,\ldots,y_K$ be
real insertion parameters.  Assume that both extended words
$((+,y_1),\ldots,(+,y_K),w)$ and
$((-,y_1+1),\ldots,(-,y_K+1),w)$ satisfy the original common
cut and separation greater than two.  Keep the source order:
lists are processed tail first.  Write these extended words as
$w_+$ and $w_-$.  Their forms are denoted by $\beta_{w_+}$ and
$\beta_{w_-}$, and $\beta_w^{[K]}$ is the componentwise extension
of the old form to the $K$ new two-dimensional factors.
Set
\[
 J=\begin{pmatrix}0&1\\-1&0\end{pmatrix},\qquad
 J_K=J^{\otimes K},\qquad
 \mathcal B_w(x,y)=\beta_w^{[K]}(J_Kx,y).
\]
Thus one new factor contributes
$\beta_w(x_1,y_0)-\beta_w(x_0,y_1)$, exactly the original
$\mathrm{pairJ}$ convention.

The determinant retained from the original construction is
\begin{equation}\label{mixed-cauchy-original-q}
 \begin{split}
 q_0(Y)&=-TT'\prod_{i=1}^k
              (Y+\tfrac12+u_i)(Y+\tfrac12+u'_i),\\
 q_w(Y)&=q_0(Y)
   \prod_{(+,a)\in w}\frac{Y-a+1}{Y-a}
   \prod_{(-,a)\in w}\frac{a-Y}{a-Y-1}.
 \end{split}
\end{equation}
The common-cut sign calculation and the stated separations give
$q_w(y_i)>0$.  No rational correction is discarded
\cite{mixed-cauchy-source,mixed-cauchy-established}.

\begin{theorem}\label{mixed-cauchy-actual}
For every original level and all compact smooth coefficient vectors
with positive Gram cutoff, or all coefficient vectors in the smooth
class with arbitrarily fast decay of every derivative and the same
cutoff,
\begin{equation}\label{mixed-cauchy-exact}
 |\mathcal B_w(x,y)|^2\prod_{i=1}^Kq_w(y_i)
           \le\beta_{w_+}(x,x)\,\beta_{w_-}(y,y).
\end{equation}
The quadratic values on the right are real and nonnegative.
This supplies the original \texttt{SizeStd} and \texttt{SizeStdCpt}
interfaces, with their original measures and normalizations.
\end{theorem}

\begin{proof}
We work on the Jacquet images of the stated compact or rapidly
decreasing group coefficient class.  The established intertwining
and preservation theorem makes this image invariant under every
operator polynomial used below, including $A_+$, $A_-$ and $C$.
Positivity is needed only on this common invariant image.
The existing contracted-matrix calculation gives the RTT relations
and both adjacent determinant identities for the current matrix
$Q_w(Y)$:
\begin{equation}\label{mixed-cauchy-two-det}
 Q_w(Y)JQ_w(Y+1)^t=q_w(Y)J,
 \qquad Q_w(Y+1)^tJQ_w(Y)=q_w(Y)J.
\end{equation}
All entries act in the old auxiliary factors and commute with the
new ones.  These identities retain their validity after every
original plus or minus update.  Indeed the fresh-factor matrices
$((Y-a)I-\Pi)/(Y-a-1)$ and
$((a-Y)I-\Omega)/(a-Y-1)$ satisfy RTT, and their determinants
in \eqref{mixed-cauchy-two-det} are respectively
$(Y-a+1)/(Y-a)$ and $(a-Y)/(a-Y-1)$.
Their coefficient algebras commute with the previous one, so the
two-dimensional alternating-tensor determinant is multiplicative.
Induction gives precisely \eqref{mixed-cauchy-original-q}.
Equivalently this follows by inserting the two displayed fresh
matrices in the two identities and multiplying their four entries.

For the following calculation relabel the new parameters in
chronological order, reversing the list if necessary.  The relabeling
only permutes fresh tensor factors and leaves their tensor $J_K$
and the product of determinants unchanged.  Put $Q_i=Q_w(y_i)$,
$Q_i^+=Q_w(y_i+1)$ in the $i$th fresh factor, and let
\[
 R_{ij}(d)=\frac{dI-\Pi_{ij}}{d-1},\qquad
 B_+=\prod_{i=2}^K\prod_{j=1}^{i-1}R_{ij}(y_i-y_j),\qquad
 B_-=\prod_{i=2}^K\prod_{j=1}^{i-1}R_{ij}(y_j-y_i).
\]
Products use the displayed increasing order.  If $A_+$ and $A_-$
are the additional residue operators beyond the old residue, their
exact recursions give
\begin{equation}\label{mixed-cauchy-fused-factors}
 A_+=Q_1\cdots Q_K B_+,
 \qquad A_-=(Q_1^+)^t\cdots(Q_K^+)^t B_-.
\end{equation}
For the plus word, move each scalar fresh-factor $R$ matrix past
the later $Q_i$ whose indices it does not meet.  For the minus word,
the successive contraction matrices multiply on the left of the
main matrix.  Partial transposition reverses that order and changes
each contraction into its permutation matrix.  Its argument is
$(y_j+1)-(y_i+1)=y_j-y_i$, proving the second formula.  The old
operator coefficients are not reordered in either calculation.

We record the two permutation identities used to cancel these
factors.  Direct expansion using the permutation relations gives
\[
 R_{ij}(d)R_{ij}(-d)=I
\]
and the Yang--Baxter identity on three distinct indices, whose
arguments are their pairwise differences.  Sliding the last row
of factors past the triangular product on the preceding indices
reverses that row; induction therefore gives
\begin{equation}\label{mixed-cauchy-r-inverse}
                            B_+B_-=B_-B_+=I.
\end{equation}
Explicitly, the induction step moves
$R_{K1}(y_K-y_1)\cdots R_{K,K-1}(y_K-y_{K-1})$
past the preceding triangular product by one three-index identity
for each crossed pair.  Its reversed product then cancels the
corresponding negative-argument row by the two-index identity.
RTT, moving one new index past each earlier one, also gives
\begin{equation}\label{mixed-cauchy-r-reversal}
        Q_1\cdots Q_K B_+=B_+ Q_K\cdots Q_1.
\end{equation}
Every denominator in these operations is nonzero under the
original separation hypotheses.

Set $q_* =\prod_iq_w(y_i)>0$.  The two identities
\eqref{mixed-cauchy-two-det} give algebraic two-sided inverses
of each $Q_i$ on the common coefficient class.  They are finite
operator matrices divided only by the nonzero scalar $q_w(y_i)$;
no closed-operator invertibility or norm bound is assumed.
As $J_K$ commutes with each permutation factor,
\[
 J_K A_-J_K^{-1}=q_*Q_1^{-1}\cdots Q_K^{-1}B_-.
\]
Now \eqref{mixed-cauchy-r-reversal} and
\eqref{mixed-cauchy-r-inverse} imply the exact fused identity
\begin{equation}\label{mixed-cauchy-fused-inverse}
       A_+C=CA_+=q_*I,
       \qquad C=J_K^{-1}A_-J_K.
\end{equation}
The two conjugations by $J_K$ agree because
$J_K^2=(-1)^KI$.

Write $\langle\cdot,\cdot\rangle_w$ for the old radial
residue form extended componentwise to the fresh factors.
The established positivity theorem for the original admissible
words says that
\[
 a(x,y)=\langle A_+x,y\rangle_w,
 \qquad b(x,y)=\langle A_-x,y\rangle_w
\]
are positive semidefinite Hermitian forms.  Thus $A_+$ and $C$
are self-adjoint relative to $\langle\cdot,\cdot\rangle_w$,
in the sense of these forms; the old form may be degenerate.
Apply Cauchy--Schwarz for $a$ to $x$ and $Cz$.  The two identities
in \eqref{mixed-cauchy-fused-inverse} give
\[
 q_*^2|\langle x,z\rangle_w|^2
       =|a(x,Cz)|^2
       \le a(x,x)a(Cz,Cz)
       =q_*a(x,x)\langle Cz,z\rangle_w.
\]
Take $z=J_K^{-1}y$ and divide by the positive scalar $q_*$.
The rightmost form becomes $b(y,y)$, and
$\langle x,J_K^{-1}y\rangle_w=\langle J_Kx,y\rangle_w$.
This proves \eqref{mixed-cauchy-exact} for radial forms.  It
also covers every zero-norm case; the old residue operator has
never been inverted.  For $K=0$ the assertion is ordinary
Cauchy--Schwarz for the old positive form.

Finally the established Jacquet compatibility theorem identifies
each actual form with its radial form times the same fixed
positive Haar constant.  The constant occurs squared on both
sides of \eqref{mixed-cauchy-exact} and cancels.  Its independence
of the word is essential here.  All operators in the proof are
finite expressions in the recovered contracted entries with
scalar spectral denominators.  They preserve the compact or
rapidly decreasing coefficient classes with their positive Gram
cutoff.  The known integrated positivity and compatibility theorems
therefore apply on the same domain throughout, and give the claimed
inequality for the original $\mathrm{formLev}$ and $\mathrm{pairJ}$.
\end{proof}

Together with the exact wave isometry, this supplies the two
actual form inequalities used by the already established finite
interpolation argument.  Identification of the partial-head
Fourier integrals, including the original root quotients at
nontrivial block strata, is still required before asserting the
complete kernel estimate.
